\begin{abstract}
  This paper presents a critical case study of the system built for the 1st Hackathon at IFPR Campus Pinhais, developed by the authors in a Software Engineering I course. Using the repository, documentation, CI pipeline, security reviews and a SonarQube analysis, we examine process, configuration management, quality, testing and maintenance. The Laravel and React system reached 93.5\% test coverage, zero bugs, vulnerabilities and code smells, and 1.4\% duplication. The main weaknesses are authorship and review concentrated in one person, late front-end tests and static analysis, and no Sonar in the CI. We close with recommendations.
\end{abstract}

\begin{resumo}
  Este artigo apresenta um estudo de caso crítico do sistema construído para o 1º Hackathon do IFPR Campus Pinhais, desenvolvido pelos autores na disciplina de Engenharia de Software I. Com base no repositório, na documentação, no \textit{pipeline} de CI, em revisões de segurança e em uma análise com SonarQube, examinamos processo, gerência de configuração, qualidade, testes e manutenção. O sistema em Laravel e React alcançou 93,5\% de cobertura, zero \textit{bugs}, vulnerabilidades e \textit{code smells} e 1,4\% de duplicação. As principais fragilidades são autoria e revisão concentradas em uma pessoa, testes de \textit{front-end} e análise estática tardios e a ausência do Sonar no CI. Concluímos com recomendações.
\end{resumo}

\textbf{Palavras-chave:} engenharia de software, estudo de caso, SonarQube, integração contínua, Laravel, hackathon.
